% !PRINTED 540
\weberpage{540}

%<<CONT
when $\omega$ runs through a complete system of roots of forms with odd
outer coefficients, with the use of the formula:
\[
f_1(\omega')f_2(\omega')f(\omega') = \sqrt{2},
\]
% [sic] The source prints -m' here but defines h' for -m immediately below (23).
if $h'$ is the class number belonging to the discriminant $-m'$:
% [sic] The middle numerator in (23) is printed 2^h; the product has h' pairs, so consistency with the last equality requires 2^{h'}.
\weq{23}{\Pi f(\omega) = \Pi f_1(\omega_1)f_2(\omega_2) = \frac{2^h}{\Pi f_1(\omega')f_2(\omega')} = \sqrt{2}^{\,h'} \Pi f(\omega'),}
in which $h$, $h'$ denote the class numbers for the discriminants $-4m$,
$-m$. We then have (§ 123):
\[
h = 2h'.
\]

If therefore, as above, $\tau$, $\tau'$ are so determined that
\[
\frac{f(\omega)}{2^\tau}, \quad \frac{f(\omega')}{2^{\tau'}}
\]
become units, then
\[
\Pi f(\omega) = 2^{h\tau}, \quad \Pi f(\omega') = 2^{h'\tau'},
\]
and from this, by (23):
\weq{24}{\tau = \frac{\tau'}{2} + \frac{1}{4},}
a formula which holds also for $m' = 1$, where $h' = h$ and
the two values $2\,\omega'$ and $\omega' : 2$ are equivalent.

By repeated application of this formula there results, when
\[
m = 4^\lambda m'
\]
holds, for an arbitrary positive $\lambda$:
\[
\tau = \frac{\tau'}{2^\lambda} + \frac{1}{2} - \frac{1}{2^{\lambda+1}}.
\]

If we summarise what has here been proved, we can
say that the following quantities are \gesperrt{algebraic units}:
\[
\begin{aligned}
\frac{f(\omega)}{\sqrt[4]{2}}, & \quad m \equiv 1, 2 \pmod 4, \\
\frac{f(\omega)}{\sqrt[3]{2}}, & \quad m \equiv 3 \pmod 8, \\
\frac{f(\omega)}{\sqrt{2}}, & \quad m \equiv 7 \pmod 8,
\end{aligned}
\]
%CONT>>
